\documentclass[10pt, letterpaper]{article}

% Packages:
\usepackage[
    ignoreheadfoot, % set margins without considering header and footer
    top=2 cm, % seperation between body and page edge from the top
    bottom=2 cm, % seperation between body and page edge from the bottom
    left=2 cm, % seperation between body and page edge from the left
    right=2 cm, % seperation between body and page edge from the right
    footskip=1.0 cm, % seperation between body and footer
    % showframe % for debugging 
]{geometry} % for adjusting page geometry
\usepackage{titlesec} % for customizing section titles
\usepackage{tabularx} % for making tables with fixed width columns
\usepackage{array} % tabularx requires this
\usepackage[dvipsnames]{xcolor} % for coloring text
\definecolor{primaryColor}{RGB}{0, 79, 144} % define primary color
\usepackage{enumitem} % for customizing lists
\usepackage{fontawesome5} % for using icons
\usepackage{amsmath} % for math
\usepackage[
    pdftitle={Scott Moon's CV},
    pdfauthor={Scott Moon},
    pdfcreator={LaTeX with RenderCV},
    colorlinks=true,
    urlcolor=primaryColor
]{hyperref} % for links, metadata and bookmarks
\usepackage[pscoord]{eso-pic} % for floating text on the page
\usepackage{calc} % for calculating lengths
\usepackage{bookmark} % for bookmarks
\usepackage{lastpage} % for getting the total number of pages
\usepackage{changepage} % for one column entries (adjustwidth environment)
\usepackage{paracol} % for two and three column entries
\usepackage{ifthen} % for conditional statements
\usepackage{needspace} % for avoiding page brake right after the section title
\usepackage{iftex} % check if engine is pdflatex, xetex or luatex

% Ensure that generate pdf is machine readable/ATS parsable:
\ifPDFTeX
    \input{glyphtounicode}
    \pdfgentounicode=1
    % \usepackage[T1]{fontenc} % this breaks sb2nov
    \usepackage[utf8]{inputenc}
    \usepackage{lmodern}
\fi



% Some settings:
\AtBeginEnvironment{adjustwidth}{\partopsep0pt} % remove space before adjustwidth environment
\pagestyle{empty} % no header or footer
\setcounter{secnumdepth}{0} % no section numbering
\setlength{\parindent}{0pt} % no indentation
\setlength{\topskip}{0pt} % no top skip
\setlength{\columnsep}{0cm} % set column seperation
\makeatletter
\let\ps@customFooterStyle\ps@plain % Copy the plain style to customFooterStyle
\patchcmd{\ps@customFooterStyle}{\thepage}{
    \color{gray}\textit{\small Scott Moon - Page \thepage{} of \pageref*{LastPage}}
}{}{} % replace number by desired string
\makeatother
\pagestyle{customFooterStyle}

\titleformat{\section}{\needspace{4\baselineskip}\bfseries\large}{}{0pt}{}[\vspace{1pt}\titlerule]

\titlespacing{\section}{
    % left space:
    -1pt
}{
    % top space:
    0.3 cm
}{
    % bottom space:
    0.2 cm
} % section title spacing

\renewcommand\labelitemi{$\circ$} % custom bullet points
\newenvironment{highlights}{
    \begin{itemize}[
        topsep=0.10 cm,
        parsep=0.10 cm,
        partopsep=0pt,
        itemsep=0pt,
        leftmargin=0.4 cm + 10pt
    ]
}{
    \end{itemize}
} % new environment for highlights

\newenvironment{highlightsforbulletentries}{
    \begin{itemize}[
        topsep=0.10 cm,
        parsep=0.10 cm,
        partopsep=0pt,
        itemsep=0pt,
        leftmargin=10pt
    ]
}{
    \end{itemize}
} % new environment for highlights for bullet entries


\newenvironment{onecolentry}{
    \begin{adjustwidth}{
        0.2 cm + 0.00001 cm
    }{
        0.2 cm + 0.00001 cm
    }
}{
    \end{adjustwidth}
} % new environment for one column entries

\newenvironment{twocolentry}[2][]{
    \onecolentry
    \def\secondColumn{#2}
    \setcolumnwidth{\fill, 4.5 cm}
    \begin{paracol}{2}
}{
    \switchcolumn \raggedleft \secondColumn
    \end{paracol}
    \endonecolentry
} % new environment for two column entries

\newenvironment{header}{
    \setlength{\topsep}{0pt}\par\kern\topsep\centering\linespread{1.5}
}{
    \par\kern\topsep
} % new environment for the header

\newcommand{\placelastupdatedtext}{% \placetextbox{<horizontal pos>}{<vertical pos>}{<stuff>}
  \AddToShipoutPictureFG*{% Add <stuff> to current page foreground
    \put(
        \LenToUnit{\paperwidth-2 cm-0.2 cm+0.05cm},
        \LenToUnit{\paperheight-1.0 cm}
    ){\vtop{{\null}\makebox[0pt][c]{
        \small\color{gray}\textit{Last updated in December 2024}\hspace{\widthof{Last updated in December 2024}}
    }}}%
  }%
}%

% save the original href command in a new command:
\let\hrefWithoutArrow\href

% new command for external links:
\renewcommand{\href}[2]{\hrefWithoutArrow{#1}{\ifthenelse{\equal{#2}{}}{ }{#2 }\raisebox{.15ex}{\footnotesize \faExternalLink*}}}


\begin{document}
    \newcommand{\AND}{\unskip
        \cleaders\copy\ANDbox\hskip\wd\ANDbox
        \ignorespaces
    }
    \newsavebox\ANDbox
    \sbox\ANDbox{}

    \placelastupdatedtext
    \begin{header}
        \textbf{\fontsize{24 pt}{24 pt}\selectfont Scott Moon}

        \vspace{0.3 cm}

        \normalsize
        \mbox{{\color{black}\footnotesize\faMapMarker*}\hspace*{0.13cm}Powder Springs, GA}%
        \kern 0.25 cm%
        \AND%
        \kern 0.25 cm%
        \mbox{\hrefWithoutArrow{mailto:scottqmoon@gmail.com}{\color{black}{\footnotesize\faEnvelope[regular]}\hspace*{0.13cm}scottqmoon@gmail.com}}%
    \end{header}

    \vspace{0.3 cm - 0.3 cm}


    \section{Research Interests}



        
        \begin{onecolentry}
            Self-organizing systems, order from chaos, complex behavior following from simple rules
        \end{onecolentry}

        \vspace{0.2 cm}

        \begin{onecolentry}
            Dynamical systems and stochastic processes
        \end{onecolentry}

        \vspace{0.2 cm}

        \begin{onecolentry}
            Mathematical biology
        \end{onecolentry}


    
    \section{Education}



        
        \begin{twocolentry}{
        \textit{West Chester, PA}    
            
        \textit{Aug 2021 – Dec 2023}}
            \textbf{West Chester University}

            \textit{MS in Applied and Computational Mathematics}
        \end{twocolentry}

        \vspace{0.10 cm}
        \begin{onecolentry}
            \begin{highlights}
                \item 4.0 GPA
            \end{highlights}
        \end{onecolentry}


        \vspace{0.2 cm}

        \begin{twocolentry}{
        \textit{Erie, PA}    
            
        \textit{Aug 2014 – July 2017}}
            \textbf{Penn State}

            \textit{BS in Mathematics}
        \end{twocolentry}




    
    \section{Presentations}



        
        \begin{onecolentry}
            ``\textit{Mayfly Mathematics: How Mayflies Synchronize}," MAA EPaDel Section Meeting, Media, PA, April 2023
        \end{onecolentry}

        \vspace{0.2 cm}

        \begin{onecolentry}
            ``\textit{Cardiac Action Potential: The Luo-Rudy Model}," WCU Research Competition, West Chester, PA, April 2022
        \end{onecolentry}


    
    \section{Research Experience}



        
        \begin{twocolentry}{
            
            
        \textit{Spring 2022 - Fall 2023}}
            \textbf{Research Assistant for Dr. Allison Kolpas}
        \end{twocolentry}

        \vspace{0.10 cm}
        \begin{onecolentry}
            \begin{highlights}
                \item Studied dynammical systems representing mayfly emergence times, looking for conditions for synchrony (stable periodic points)
                \item Developed a model for predicting the final weight of mayflies, used as a proxy for fecundity in population study
                \item Programmed an agent based model to simulate synchrony and population dynamics
            \end{highlights}
        \end{onecolentry}


        \vspace{0.2 cm}

        \begin{twocolentry}{
            
            
        \textit{Fall 2022}}
            \textbf{Research Aid for Dr. Frank Fish}
        \end{twocolentry}

        \vspace{0.10 cm}
        \begin{onecolentry}
            \begin{highlights}
                \item Created a program for measuring turning radius of swimming turtles from video data
            \end{highlights}
        \end{onecolentry}


        \vspace{0.2 cm}

        \begin{twocolentry}{
            
            
        \textit{Fall 2016}}
            \textbf{Undergraduate research project with Dr. Joseph Previte}
        \end{twocolentry}

        \vspace{0.10 cm}
        \begin{onecolentry}
            \begin{highlights}
                \item Studied properties of certain matrices arising in finite calculus
            \end{highlights}
        \end{onecolentry}



    
    \section{Professional Experience}



        
        \begin{twocolentry}{
            
            
        \textit{2024 – present}}
            \textbf{Senior Data Analyst}
            
            \textit{Charles Schwab}
        \end{twocolentry}

        \vspace{0.10 cm}
        \begin{onecolentry}
            \begin{highlights}
                \item Statistical analysis for classifying customer web and trade behavior
            \end{highlights}
        \end{onecolentry}


        \vspace{0.2 cm}

        \begin{twocolentry}{
            
            
        \textit{2022 – 2023}}
            \textbf{Data Engineer}
            
            \textit{Hagerty}
        \end{twocolentry}

        \vspace{0.10 cm}
        \begin{onecolentry}
            \begin{highlights}
                \item Developed Airflow ETL pipelines in Python, used Snowflake for data warehousing, and used DBT for view creation
                \item Used Terraform for infrastructure-as-code for AWS resources
                \item Created custom Python library for automated schema management, metadata, and dynamic SQL queries
                \item Wrote technical documents for architecture and development proposals
            \end{highlights}
        \end{onecolentry}


        \vspace{0.2 cm}

        \begin{twocolentry}{
            
            
        \textit{2020 – 2022}}
            \textbf{Data Engineer}
            
            \textit{Vanguard}
        \end{twocolentry}

        \vspace{0.10 cm}
        \begin{onecolentry}
            \begin{highlights}
                \item Used PySpark for ETL of unstructured and structured data; running on AWS EMR cluster
                \item Set up and managed backend for ThoughtSpot BI tool running on AWS EC2
                \item Implemented infrastructure-as-code using CloudFormation and CDK in Python
                \item Created data governance system with Boto3, IAM, and Glue Data Catalog
                \item Created custom data quality system to check completeness, accuracy, and timeliness
            \end{highlights}
        \end{onecolentry}


        \vspace{0.2 cm}

        \begin{twocolentry}{
            
            
        \textit{2017 – 2020}}
            \textbf{Modeling Analyst and Controls Engineer}
            
            \textit{Enersys}
        \end{twocolentry}

        \vspace{0.10 cm}
        \begin{onecolentry}
            \begin{highlights}
                \item Designed the ESS (Energy Storage System) architecture and wrote the specifications
                \item Created an ESS simulation to model cost savings, including a web applet using React and Django
                \item Implemented the ESS control program in Python running on an embedded Linux system
                \item Architected PostgreSQL database for storing system data and managing live communication
                \item Created a Node web application and REST API for monitoring and controlling the ESS
                \item Led an FMEA and analyzed 6 months of live testing data with R
            \end{highlights}
        \end{onecolentry}



    
    \section{Grants and Awards}

    \begin{onecolentry}
        \begin{highlightsforbulletentries}


        \item Viorel Nitica Scholarship, 2022

        \item 1st place in West Chester Research Competition, 2022

        \item Undergraduate Research Award, 2016


        \end{highlightsforbulletentries}
    \end{onecolentry}

    \section{Extracirricular Activies}



        
        \begin{onecolentry}
            \textbf{Math Club}, \textit{President} (2016-2017) - organized faculty talks and fundraising events
        \end{onecolentry}

        \vspace{0.2 cm}

        \begin{onecolentry}
            \textbf{International Student Organization}, \textit{Program Director} (2016-2017) - organized community building events, celebratory dinners, and fundraising events
        \end{onecolentry}

        \vspace{0.2 cm}

        \begin{onecolentry}
            \textbf{Student Government Association}, \textit{Budget Committee Member} (2015-2017)
        \end{onecolentry}

        \vspace{0.2 cm}

        \begin{onecolentry}
            \textbf{Musician's Club}, \textit{Vice President} (2015-2017) - organized student concerts and coordinated musical entertainment at other student events
        \end{onecolentry}

        \vspace{0.2 cm}

        \begin{onecolentry}
            \textbf{Math Tutor} - drop-in mathematics tutor at PSU (2015-2016) and WCU (2021-2022)
        \end{onecolentry}


    
    \section{Projects}



        
        \begin{onecolentry}
            \textbf{Turing Morphogenesis}
        \end{onecolentry}

        \vspace{0.10 cm}
        \begin{onecolentry}
            \begin{highlights}
                \item Exposition on Alan Turing's original 1952 reaction-diffusion paper from a modern perspective, including a bifurcation analysis and 2D numerical simulation using parallelism.
            \end{highlights}
        \end{onecolentry}


        \vspace{0.2 cm}

        \begin{onecolentry}
            \textbf{Cardiac Action Potential: The Luo-Rudy Model}
        \end{onecolentry}

        \vspace{0.10 cm}
        \begin{onecolentry}
            \begin{highlights}
                \item A simulation and analysis of the Luo-Rudy model (a Hodgkin-Huxley type model for heart cells) propagating through a cardiac fiber using the cable equation
            \end{highlights}
        \end{onecolentry}


        \vspace{0.2 cm}

        \begin{onecolentry}
            \textbf{Iterated Games and Markov Chains}
        \end{onecolentry}

        \vspace{0.10 cm}
        \begin{onecolentry}
            \begin{highlights}
                \item An analysis of reactive strategies for the Prisoner's Dilemma using Markov Chains and evolutionary game theory
            \end{highlights}
        \end{onecolentry}


        \vspace{0.2 cm}

        \begin{onecolentry}
            \textbf{Walls, an abstract strategy game}
        \end{onecolentry}

        \vspace{0.10 cm}
        \begin{onecolentry}
            \begin{highlights}
                \item A board game inspired by Theseus and Minotaur, in which both players can move walls and change the maze; wrote a C\# implementation and built a physical wooden board and pieces
            \end{highlights}
        \end{onecolentry}



    
    \section{Technologies}



        
        \begin{onecolentry}
            \textbf{Languages:} Python, MATLAB, R, Fortran, JavaScript, C\#, C++ C, SQL, HTML, CSS
        \end{onecolentry}

        \vspace{0.2 cm}

        \begin{onecolentry}
            \textbf{Frameworks and Platforms:} Pyspark, AWS, Hadoop, Airflow, Terraform, REST, Django, GraphQL, React
        \end{onecolentry}


    

\end{document}